\section{Marginal CARMA Calibration}\label{sec:marginal-carma-calibration}

\subsection{Calibration Framework}\label{ssec:carma-calibration-framework}

Each deseasonalized marginal factor is modeled as a stationary Levy-driven
CARMA process in state-space form \citep{BrockwellLindner2024Ch18}. For a
CARMA\((p,q)\),
\begin{equation}
  dG_t=A G_t\,dt+e_p\,dL_t,
  \qquad
  X_t=b^\top G_t ,
  \label{eq:carma-state-space}
\end{equation}
where \(L\) is a one-dimensional Levy driver and \(A\) is the companion matrix
of
\begin{equation}
  a(z)=z^p+a_1z^{p-1}+\cdots+a_p .
  \label{eq:carma-ar-polynomial}
\end{equation}
The roots of \(a\) are required to have negative real parts. The moving-average
vector is written as \(b=(b_0,\ldots,b_q,0,\ldots,0)^\top\), with the last
non-zero coefficient normalized to one. Throughout the paper we use the full
moving-average convention \(q=p-1\), unless stated otherwise.

The order is initialized from the empirical autocorrelation function. A real
root contributes a monotone exponential decay, while a complex conjugate pair
contributes a damped oscillation. We therefore fit
\begin{equation}
  \rho_{\mathrm{ms}}(h)
  =
  \sum_{r=1}^{R} w_r e^{-\kappa_r h}
  +
  \sum_{\ell=1}^{C} w_{o,\ell} e^{-\kappa_{o,\ell}h}
  \cos(\omega_\ell h),
  \qquad
  \sum_r w_r+\sum_\ell w_{o,\ell}=1 ,
  \label{eq:carma-multiscale-acf}
\end{equation}
with non-negative weights. In the applications below there is one daily
oscillatory pair, with \(\omega=2\pi/24\), and the fitted ACF is spectrally
factorized to obtain a stable initial CARMA model.

Starting from this initialization, the final marginal parameters are estimated
by exact Gaussian prediction-error QMLE, following the Kalman likelihood
construction for sampled CARMA processes \citep{BrockwellLindner2024Ch19}. For
given CARMA coefficients, the Kalman recursion gives one-step prediction errors
\(e_i\) and innovation variances \(r_i\) under unit Levy variance. Since the
Levy drift enters as a constant mean shift, the adjusted innovation is
\[
  \varepsilon_i(m)=e_i-mc_i .
\]
The drift rate and variance rate are profiled as
\begin{equation}
  \widehat m
  =
  \frac{\sum_i e_i c_i/r_i}{\sum_i c_i^2/r_i},
  \qquad
  \widehat \nu^2
  =
  \frac{1}{n}\sum_i
  \frac{\varepsilon_i(\widehat m)^2}{r_i}.
  \label{eq:carma-qmle-profile}
\end{equation}
The optimized criterion is the reduced Gaussian prediction-error likelihood,
augmented when needed by a soft ACF penalty,
\begin{equation}
  \mathcal J(\theta)
  =
  \log(\widehat\nu^2)
  +
  \frac{1}{n}\sum_i \log r_i
  +
  \lambda_{\mathrm{acf}}
  \frac{1}{|H|}
  \sum_{h\in H}
  \omega_h
  \left(\rho_\theta(h)-\widehat\rho(h)\right)^2 .
  \label{eq:carma-qmle-objective}
\end{equation}
The penalty stabilizes long-lag dependence without replacing the likelihood
criterion. Its value is reported separately for each marginal.

Once the CARMA parameters are fixed, the hourly Levy increments are recovered
from the smoothed state path, as in the QMLE framework of
\citet{BrockwellLindner2024Ch19}. Let \(\widehat G_t\) denote the RTS-smoothed
state, including the profiled drift correction \(\widehat m(-A)^{-1}e_p\).
Choose a real root \(\lambda_r\) of \(a\), taken in practice as the slowest real
root, and let \(E\) be the matrix of companion eigenvectors. The corresponding
modal coordinate satisfies a scalar driven equation, giving
\begin{equation}
  \widehat Z_r(t)=b(\lambda_r)e_r^\top E^{-1}\widehat G_t,\qquad
  d\widehat Z_r(t)=\lambda_r\widehat Z_r(t)\,dt+\alpha_r\,dL_t,
  \qquad
  \alpha_r=\frac{b(\lambda_r)}{a'(\lambda_r)},
  \label{eq:levy-modal-recovery}
\end{equation}
and therefore
\begin{equation}
  \widehat L_t-\widehat L_0
  =
  \frac{\widehat Z_r(t)-\widehat Z_r(0)
  -\lambda_r\int_0^t \widehat Z_r(s)\,ds}{\alpha_r}.
  \label{eq:levy-driver-recovery}
\end{equation}
On the hourly grid the integral is evaluated numerically, and the empirical
driver sample used below is
\(\Delta\widehat L_i=\widehat L_{t_i}-\widehat L_{t_{i-1}}\).

\subsection{German Spot Prices}\label{ssec:german-price-carma-calibration}

The deseasonalized German log-price factor \(\widehat X^P\) is calibrated with
the procedure of \cref{ssec:carma-calibration-framework}. To represent a
decaying ACF with damped 24-hour oscillations, the initialization combines one
daily oscillatory pair with three real half-lives. The first real scale is
expected to be hourly, the second daily, and the third is fixed at 120 d during
initialization to impose a slow component for option-pricing horizons.
\begin{equation}
  \rho_{\mathrm{ms}}(h)
  =
  \sum_{i=1}^{3} w_i \exp(-\kappa_i h)
  +
  w_4 \exp(-\kappa_4 h)
  \cos\left(\frac{2\pi h}{24}\right),
  \qquad
  \sum_{i=1}^{4}w_i=1,\quad w_i\ge 0,
  \label{eq:price-multiscale-acf}
\end{equation}
Here \(h\) is measured in hours, and the oscillatory frequency is fixed at
\(2\pi/24\). The additional long real half-life is specific to the price
model. It is introduced only at initialization to keep slow price persistence
available for pricing, while the temperature, solar, and wind marginals below
do not require such an imposed long mode. The fitted decomposition has three
real roots and one complex conjugate pair, which gives the autoregressive order
\[
  p = 3 + 2 = 5 .
\]
We use the full moving-average order \(q=p-1=4\). Spectral factorization of the
fitted ACF then gives the initial CARMA(5,4) model.

\FloatBarrier

\begin{table}[!htbp]
  \centering
  \caption{German price CARMA(5,4) half-lives before and after penalized QMLE.}
  \label{tab:price-carma-half-lives}
  \begin{tabular}{@{}lll@{}}
    \toprule
    Root & Half-life after initialization & Half-life after QMLE \\
    \midrule
    \(\lambda_1\) & \(2.63\) h & \(4.52\) h \\
    \(\lambda_2\) & \(1.20\) d & \(1.50\) d \\
    \(\lambda_o\) & \(7.21\) d & \(7.21\) d \\
    \(\lambda_3\) & \(120\) d & \(82.4\) d \\
    \bottomrule
  \end{tabular}
\end{table}

\FloatBarrier

The QMLE step uses \(\lambda_{\mathrm{acf}}=10\). The daily pair remains fixed,
whereas the three real half-lives, including the slow one, and the
moving-average coefficients are re-estimated. The final log-likelihood is
\(77709\). As shown in
\cref{fig:price-spot-acf-calibration}, the ACF penalty keeps the final CARMA
autocorrelation close to the multiscale initialization while allowing the
likelihood to improve.

\begin{figure}[!htbp]
  \centering
  \includegraphics[width=0.78\textwidth]{calibration/price_spot_acf_calibration_300.pdf}
  \caption{German price ACF calibration over 300 hourly lags, comparing the
  empirical autocorrelation, the multiscale initialization, and the final
  CARMA(5,4) autocorrelation after QMLE.}
  \label{fig:price-spot-acf-calibration}
\end{figure}

\begin{table}[!htbp]
  \centering
  \caption{Gaussian log-likelihood comparison on the same deseasonalized German
  log-price residuals.}
  \label{tab:price-loglik-comparison}
  \begin{tabular}{@{}lr@{}}
    \toprule
    Model & Log-likelihood \\
    \midrule
    CARMA(5,4), ACF initialization & \(77562\) \\
    CARMA(5,4), QMLE & \(77709\) \\
    ARMA(1,0) & \(77454\) \\
    ARMA(5,4) & \(78661\) \\
    \bottomrule
  \end{tabular}
\end{table}

\Cref{tab:price-loglik-comparison} shows that the penalized QMLE improves the
Gaussian log-likelihood relative to the ACF initialization, although it remains
below the unconstrained ARMA(5,4) benchmark. Pure QMLE can increase this
likelihood further, as discussed below, but only by shortening the slow price
component too much.

\begin{figure}[!htbp]
  \centering
  \includegraphics[width=0.70\textwidth]{calibration/price_spot_driver_distribution_fit.pdf}
  \caption{Recovered hourly Levy driver increments for the German price CARMA
  model, with Gaussian and NIG distribution fits.}
  \label{fig:price-spot-driver-distribution}
\end{figure}

After calibration, the driver increments recovered by
\cref{eq:levy-driver-recovery} are fitted with a Gaussian law and with a normal
inverse Gaussian law. The NIG specification is preferred according to the AIC
criterion, consistently with \cref{fig:price-spot-driver-distribution}. The NIG
density better matches the asymmetry and tail thickness of the recovered
increments than the Gaussian benchmark.

\begin{figure}[!htbp]
  \centering
  \includegraphics[width=0.74\textwidth]{calibration/price_spot_final_price_distribution.pdf}
  \caption{Distribution check for the final German spot price in EUR/MWh,
  comparing the observed data with prices simulated from the Gaussian and NIG
  CARMA marginals.}
  \label{fig:price-spot-final-price-distribution}
\end{figure}

\Cref{fig:price-spot-final-price-distribution} maps the simulated CARMA
residuals back to EUR/MWh using the same seasonal component. The NIG-driven
model improves the distribution relative to the Gaussian-driven model, but
neither specification reproduces the bimodality of the empirical price
distribution, in particular the concentration around zero EUR/MWh. In the
European day-ahead market, hourly prices are obtained by uniform-price auctions
from the intersection of aggregated supply and demand curves. During hours with
low residual demand, renewable production and other low marginal-cost or
must-run units frequently set, or strongly constrain, the clearing price near
zero. During tighter hours, the marginal unit is usually thermal or flexible
generation with a strictly positive bid. This market-clearing mechanism creates
regime-like mass in the final price distribution and cannot be reproduced by a
smooth continuous-time CARMA marginal factor alone.

We also ran a pure exact QMLE without the ACF penalty
\((\lambda_{\mathrm{acf}}=0)\), using semi-global restarts from the same
CARMA(5,4) structure. The best run over 20 restarts increases the Gaussian
log-likelihood because the objective only rewards one-step prediction fit. It
does not penalize the loss of long-lag ACF structure.

\begin{table}[H]
  \centering
  \caption{Penalized versus pure QMLE for the German price CARMA(5,4).}
  \label{tab:price-pure-qmle-comparison}
  \begin{tabular}{@{}lcc@{}}
    \toprule
    Quantity & Retained QMLE & Pure QMLE \\
    \midrule
    ACF penalty \(\lambda_{\mathrm{acf}}\) & \(10\) & \(0\) \\
    Log-likelihood & \(77709\) & \(78808\) \\
    \(\lambda_1\) half-life & \(4.52\) h & \(0.510\) h \\
    \(\lambda_2\) half-life & \(1.50\) d & \(1.20\) h \\
    \(\lambda_o\) half-life & \(7.21\) d & \(5.20\) d \\
    \(\lambda_3\) half-life & \(82.4\) d & \(1.40\) d \\
    \bottomrule
  \end{tabular}
\end{table}

The pure QMLE is therefore better as a one-step Gaussian likelihood fit, but it
collapses the slow price factor from several months to about one day. This
choice is not appropriate for the pricing step, because the model would lose the
long-horizon persistence imposed by the ACF initialization and retained by the
penalized QMLE.

\FloatBarrier

\subsection{German Temperature}\label{ssec:german-temperature-carma-calibration}

The stochastic temperature factor \(\widehat X^T\) is the hourly Berlin
temperature residual after deterministic Fourier deseasonalization. It is
modeled with the stationary CARMA structure of \cref{eq:carma-state-space}. Its
empirical ACF is smoother than the electricity price ACF, but it still contains
a daily oscillatory component. The initialization therefore uses two real modes
and one daily pair.
\begin{equation}
  \rho_{\mathrm{ms}}^T(h)
  =
  w_f e^{-\kappa_f h}
  +
  w_s e^{-\kappa_s h}
  +
  w_d e^{-\kappa_d h}
  \cos\left(\frac{2\pi h}{24}\right),
  \qquad
  w_f+w_s+w_d=1,\quad w_f,w_s,w_d\ge 0 .
  \label{eq:temperature-multiscale-acf}
\end{equation}
The two real components represent a relatively fast decay and a slower
persistent component, while the complex conjugate pair enforces the 24-hour
oscillation. The resulting order is
\[
  p = 2 + 2 = 4,
  \qquad q=p-1=3 .
\]
The daily frequency is fixed in the ACF step and the daily AR pair is kept fixed
in QMLE. The two real AR roots and the moving-average polynomial are
re-estimated.

\FloatBarrier

\begin{table}[!htbp]
  \centering
  \caption{German temperature CARMA(4,3) half-lives before and after penalized
  QMLE.}
  \label{tab:temperature-carma-half-lives}
  \begin{tabular}{@{}lll@{}}
    \toprule
    Root & Half-life after initialization & Half-life after QMLE \\
    \midrule
    \(\lambda_1\) & \(1.06\) d & \(0.530\) h \\
    \(\lambda_2\) & \(2.42\) d & \(1.59\) d \\
    \(\lambda_o\) & \(3.03\) d & \(3.03\) d \\
    \bottomrule
  \end{tabular}
\end{table}

\FloatBarrier

\begin{figure}[!htbp]
  \centering
  \includegraphics[width=0.78\textwidth]{calibration/temperature_acf_calibration_300.pdf}
  \caption{German temperature ACF calibration over 300 hourly lags, comparing
  the empirical autocorrelation, the multiscale initialization, and the final
  CARMA(4,3) autocorrelation after QMLE.}
  \label{fig:temperature-acf-calibration}
\end{figure}

For temperature, the QMLE step uses \(\lambda_{\mathrm{acf}}=0.1\). This
penalty is lower than for prices because the Gaussian likelihood is on a
different numerical scale. The trade-off is visible in
\cref{fig:temperature-acf-calibration}, where the final ACF is less tightly
matched than in the price calibration. The weaker penalty nevertheless improves the
likelihood fit relative to the ACF initialization, with a final log-likelihood
of \(-22941\), while remaining below the ARMA(4,3) benchmark.

\begin{table}[!htbp]
  \centering
  \caption{Gaussian log-likelihood comparison on the same deseasonalized German
  temperature residuals.}
  \label{tab:temperature-loglik-comparison}
  \begin{tabular}{@{}lr@{}}
    \toprule
    Model & Log-likelihood \\
    \midrule
    CARMA(4,3), ACF initialization & \(-23760\) \\
    CARMA(4,3), QMLE & \(-22941\) \\
    ARMA(1,0) & \(-24612\) \\
    ARMA(4,3) & \(-23173\) \\
    \bottomrule
  \end{tabular}
\end{table}

After QMLE, the slow real half-life is about one and a half days, close to the
solar value and much shorter than the wind and price slow modes. Temperature is
therefore the cleanest two-scale physical marginal, with a fast adjustment, a
one-to-two-day persistence scale, and the fixed daily oscillation.

\begin{figure}[H]
  \centering
  \includegraphics[width=0.78\textwidth]{calibration/temperature_driver_distribution_fit.pdf}
  \caption{Recovered hourly Levy driver increments for the German temperature
  CARMA model. The NIG fit captures the asymmetric and heavy-tailed driver
  distribution better than the Gaussian fit.}
  \label{fig:temperature-driver-distribution}
\end{figure}

\FloatBarrier

\begin{figure}[!htbp]
  \centering
  \includegraphics[width=0.74\textwidth]{calibration/temperature_final_distribution.pdf}
  \caption{Distribution check for the final German temperature level, comparing
  the observed data with paths simulated from the Gaussian and NIG CARMA
  marginals.}
  \label{fig:temperature-distribution-checks}
\end{figure}

\FloatBarrier

\subsection{German Solar Production}\label{ssec:german-solar-carma-calibration}

The solar marginal model is fitted to the deseasonalized latent residual
\(\widehat X^Q\) defined after the clear-sky normalization and logit
transformation in \cref{ssec:solar-deseasonalization}. The physical capacity
factor is bounded and has structural nighttime zeros, so the CARMA model is not
applied directly to \(Q_t\). It is estimated in the transformed coordinate, and
simulated paths are mapped back through the inverse clear-sky transformation.

As for temperature, we use a CARMA(4,3) with two real modes and one daily
oscillatory pair. The solar ACF is more specific because it contains periodic
peaks every 24 hours, which must be preserved in calibration. The daily pair is
fixed after the ACF step, while QMLE re-estimates the two real persistence
scales and the moving-average polynomial.

\FloatBarrier

\begin{table}[!htbp]
  \centering
  \caption{German solar CARMA(4,3) half-lives before and after penalized QMLE.}
  \label{tab:solar-carma-half-lives}
  \begin{tabular}{@{}lll@{}}
    \toprule
    Root & Half-life after initialization & Half-life after QMLE \\
    \midrule
    \(\lambda_1\) & \(0.506\) h & \(0.512\) h \\
    \(\lambda_2\) & \(1.74\) d & \(1.57\) d \\
    \(\lambda_o\) & \(2.93\) d & \(2.93\) d \\
    \bottomrule
  \end{tabular}
\end{table}

\FloatBarrier

\begin{figure}[!htbp]
  \centering
  \includegraphics[width=0.78\textwidth]{calibration/solar_acf_calibration_300.pdf}
  \caption{German solar latent-residual ACF calibration over 300 hourly lags,
  comparing the empirical autocorrelation, the multiscale initialization, and
  the final CARMA(4,3) autocorrelation after QMLE.}
  \label{fig:solar-acf-calibration}
\end{figure}

For solar, the QMLE step uses \(\lambda_{\mathrm{acf}}=1000\). This larger
penalty is chosen to preserve the latent ACF structure, in particular its daily
oscillatory component. This choice is also supported by the likelihood
comparison, as the ACF initialization already has a higher log-likelihood than
the ARMA(4,3) benchmark. QMLE then refines the fit only slightly, with a final
log-likelihood of \(-84623\). This leads to only small changes in both the
estimated half-lives and the log-likelihood.

\begin{table}[!htbp]
  \centering
  \caption{Gaussian log-likelihood comparison on the same deseasonalized German
  solar latent residuals.}
  \label{tab:solar-loglik-comparison}
  \begin{tabular}{@{}lr@{}}
    \toprule
    Model & Log-likelihood \\
    \midrule
    CARMA(4,3), ACF initialization & \(-84637\) \\
    CARMA(4,3), QMLE & \(-84623\) \\
    ARMA(1,0) & \(-86462\) \\
    ARMA(4,3) & \(-84786\) \\
    \bottomrule
  \end{tabular}
\end{table}

The solar CARMA(4,3) has roots close to those obtained for temperature. One
real half-life is below one hour, another is around one and a half days, and
the daily oscillations have a damping half-life of about three days.

\begin{figure}[H]
  \centering
  \includegraphics[width=0.78\textwidth]{calibration/solar_driver_distribution_fit.pdf}
  \caption{Recovered hourly Levy driver increments for the German solar CARMA
  model. The NIG fit captures the asymmetric and heavy-tailed driver
  distribution better than the Gaussian fit.}
  \label{fig:solar-driver-distribution}
\end{figure}

\FloatBarrier

\subsection{German Wind Production}\label{ssec:german-wind-carma-calibration}

The wind marginal model is fitted to the deseasonalized logit residual
\(\widehat X^W\) defined in \cref{ssec:wind-deseasonalization}. As for solar,
the CARMA model is estimated in an unconstrained latent coordinate and simulated
paths are mapped back to capacity factors with the inverse logit transform. This
keeps the continuous-time model away from the physical bounds \(0\) and \(1\).

As for temperature and solar, we use a CARMA(4,3) with two real modes and one
daily oscillatory pair. The wind ACF is less tied to a deterministic day-night
structure than solar, but it has longer persistence. The daily pair is fixed
after the ACF step, while QMLE re-estimates the two real persistence scales and
the moving-average polynomial.

\FloatBarrier

\begin{table}[!htbp]
  \centering
  \caption{German wind CARMA(4,3) half-lives before and after penalized QMLE.}
  \label{tab:wind-carma-half-lives}
  \begin{tabular}{@{}lll@{}}
    \toprule
    Root & Half-life after initialization & Half-life after QMLE \\
    \midrule
    \(\lambda_1\) & \(16.1\) h & \(20.9\) h \\
    \(\lambda_o\) & \(3.28\) d & \(3.28\) d \\
    \(\lambda_2\) & \(7.44\) d & \(37.6\) d \\
    \bottomrule
  \end{tabular}
\end{table}

\FloatBarrier

\begin{figure}[!htbp]
  \centering
  \includegraphics[width=0.78\textwidth]{calibration/wind_acf_calibration_300.pdf}
  \caption{German wind logit-residual ACF calibration over 300 hourly lags,
  comparing the empirical autocorrelation, the multiscale initialization, and
  the final CARMA(4,3) autocorrelation after QMLE.}
  \label{fig:wind-acf-calibration}
\end{figure}

For wind, the QMLE step uses \(\lambda_{\mathrm{acf}}=0.1\), giving a final
log-likelihood of \(-6984\). The fit improves the ACF initialization but
remains below the ARMA(4,3) benchmark in log-likelihood.

\begin{table}[!htbp]
  \centering
  \caption{Gaussian log-likelihood comparison on the same deseasonalized German
  wind logit residuals.}
  \label{tab:wind-loglik-comparison}
  \begin{tabular}{@{}lr@{}}
    \toprule
    Model & Log-likelihood \\
    \midrule
    CARMA(4,3), ACF initialization & \(-7013\) \\
    CARMA(4,3), QMLE & \(-6984\) \\
    ARMA(1,0) & \(-7262\) \\
    ARMA(4,3) & \(-6853\) \\
    \bottomrule
  \end{tabular}
\end{table}

The main difference with solar and temperature is the slow real mode, whose
half-life increases to about 38 days after QMLE. It represents persistent
wind-regime variation in the latent residual.

\begin{figure}[H]
  \centering
  \includegraphics[width=0.78\textwidth]{calibration/wind_driver_distribution_fit.pdf}
  \caption{Recovered hourly Levy driver increments for the German wind CARMA
  model. The NIG fit captures the narrow center and heavy tails better than the
  Gaussian fit.}
  \label{fig:wind-driver-distribution}
\end{figure}

\FloatBarrier

\begin{figure}[!htbp]
  \centering
  \includegraphics[width=0.74\textwidth]{calibration/wind_final_distribution.pdf}
  \caption{Distribution check for the final German wind capacity factor,
  comparing the observed data with paths simulated from the Gaussian and NIG
  CARMA marginals.}
  \label{fig:wind-distribution-checks}
\end{figure}

\FloatBarrier
